\def\ChapterCode{KOM}
\chapter
[\CO{[KOM]} Kommunikation mit Patienten, Angehörigen und Mitarbeitern] 
{Kommunikation mit Patienten, Angehörigen und Mitarbeitern}
\label{chp:KOM}

\ChapterIntro
{%

}
{%
\begin{description}
  \item[Maintainer:] Sebastian Gabriel
  \item[Autoren:] Diverse
  \item[Reviewer:] Standard-Reviewprozess
  \item[Version:] {\VarDocVersion} ({\VarDocVersionDescription})
  \item[SHA1:] {\DocGitCommit}
\end{description}

{\TxTeilkapitelAASS}}

\clearpage

\section{Kommunikationsregeln}
\label{topic:kommunikationsregeln}

\dictum[Paul Watzlawick, \cite{WatzlawickMenschlKomm11}]{Man
kann nicht \E{nicht} kommunizieren.}\bigskip

Kommunikation ist die Voraussetzung für eine erfolgreiche
Zusammenarbeit in einer Gemeinschaft. Bei der Versorgung
eines Patienten gilt dieser Grundsatz ebenso, zusätzlich
finden sich hier spezielle Herausforderungen. Einerseits
besteht diese \enquote*{Gemeinschaft} aus sehr
unterschiedlichen Personen: Fachpersonal unterschiedlicher
Ausbildungsstufen einerseits, und Patienten und Angehörige,
welche zumeist über wenig Fachwissen, dafür aber oft über
ausgeprägte und mal mehr, mal weniger verständliche
Eigeninteressen verfügen andererseits. Dazu kommen mitunter
noch weitere Personen wie Polizisten, Feuerwehrleute,
Bedienstete von Verkehrsbetrieben, Dritte (Berufer,
Schaulustige, Zeugen, \ldots), usw., mit welchen
kommuniziert werden muss. 
Als wäre die Fülle von potentiellen Kommunikationspartnern
nicht schon Herausforderung genug, findet man sich häufig in einer
Situation wieder, die für die Kommunikationspartner einen
Ausnahmezustand darstellt: Krankheit, Schmerzen,
viele unbekannte Gesichter, unbekannte Geräte, eine oft
unbekannte Fachsprache, die Ungewissheit was passieren wird
{\ldots} All dem soll im Rahmen einer professionellen
Kommunikation entsprechend begegnet werden. 
Ziele der Kommunikation sind:

\begin{itemize}
  \item Ermittlung von Informationen
\begin{compactitem}
  \item Situation
  \item Anamnese (mit Patient, Fremdanamnese, \ldots)
\end{compactitem}
  \item \EE{Aufklärung} des Patienten: Jeder Patient ist vor
  Durchführung einer Untersuchung oder Maßnahme über
  diese aufzuklären.
  Das \E{Recht auf Aufklärung} ist nach der
  \E{Grundrechtecharta} der \Abk{EU} ein \EE{Grundrecht}
  (\FullPrettyRef{topic:aufklaerung}).
  
  Vor allem ist es aber eine Frage des Respekts  gegenüber
  dem Patienten ihm zu sagen, was man an und mit \E{seinem}
  Körper macht:
  \begin{compactitem}
    \item Untersuchungen
    \item Maßnahmen
    \item Weitere Vorgänge
  \end{compactitem}
  \item \EE{Psychische Betreuung} -- Eng in Verbindung mit
  der Aufklärung
  \begin{compactitem}
    \item Bildung einer Vertrauensbasis
    \item Förderung der Kooperation
  \end{compactitem}
  \item \EE{Verständigung} innerhalb des Teams
  \begin{compactitem}
    \item \enquote{Alle Teammitgleider am gleichen Stand}
    \item Absprache des weiteren Vorgehens
    \item Koordination von Untersuchungen und Maßnahmen
  \end{compactitem}
  \item \EE{Übergabe von Informationen} an weiterbehandelnde
  Stellen
\end{itemize}

\section{Kurze Kommunikationstheorie}

Es gibt Theorien zur Kommunikation wie Sand am Meer. Im
folgenden wollen wir uns auf zwei wesentliche, in der Praxis
sehr taugliche Modelle konzentrieren, welche viele der
Fehler und Fallen, die tagtäglich auf uns lauern, erklären
können. Dabei stehen zwei Aspekte im Vordergrund:
\begin{enumerate}
  \item \EE{Übermittlung}: Es gibt bei der
  Nachrichtenübermittlung einen Sender und einen Empfänger.
  \item \EE{Aussage}: Eine Nachricht enthält
  unterschiedliche Aussagen.
\end{enumerate}

\Text
{Übermittlung: Sender und Empfänger}
{%
Das Sender\,--\,Empfänger-Modell erklärt das Phänomen der 
eigentlichen Nachrichtenübermittlung \cite{Informatik3}.
Eine \EE{Nachricht} wird dabei 
\begin{enumerate}
  \item von
  einem \EE{Sender}
  \item \EE{kodiert} (Sprache, Fachsprache, Dialekt,
  Umgangssprache, Alphabet, \ldots)
  \item und über einen \EE{Übertragungsweg} (direkt mündlich,
  Sprechfunk, Datenfunk, Protokolle und Niederschriften, \ldots)
  \item zu einem \EE{Empfänger} übermittelt
  \item und dort \EE{dekodiert}.
\end{enumerate}
Die Nachricht kann \EE{durch Störungen verfälscht} werden.
Diese Störung kann auftreten:
\begin{itemize}
  \item Beim \EE{Sender}: Der Sender sagt nicht was er sagen
  will, undeutliche Sprache, Sprachschwieirgkietn, 
  falsch verstandene Fachsprache, \ldots
  \item Bei der \EE{Kodierung}: Mehrdeutigkeit, kulturelle
  Unterschiede,
  \item Beim \EE{Übertragungsweg}: Laute Umgebung,
  undeutlicher Sprechfunk, unleserliche Schrift, \ldots
  \item Beim \EE{Empfänger} und bei der \EE{Dekodierung}:
  Unkenntnis der Fachsprache oder der Umgangssprache, Schwerhörigkeit, Ablenkung, Unkenntnis
  des Zusammenhangs, eingegrenzte Wahrnehmung, \ldots
  
  Ein großes Problem dabei ist, dass der Empfänger oft
  nicht bemerkt, dass er etwas falsch verstanden hat.
\end{itemize}

\Fazit{\E{Gedacht} ist nicht gesagt, 
\E{gesagt} ist nicht gehört,
\E{gehört} ist nicht verstanden,
\E{verstanden} ist nicht gewollt.}

Ein gutes Beispiel einer Störung der Kommunikation ist das
\enquote*{Stille-Post}-Spiel, bei welchem am Ende der
Übertragungskette die ursprüngliche Nachricht zumeist völlig
entstellt ist. Um deratige Übertragungsfehler zu
\E{vermeiden}, ist es wichtig, dass der Empfänger eine
Kontroll- bzw. \EE{Rückmeldung} an den Sender gibt, um zu
kontrollieren, ob und wie er die Botschaft verstanden hat.

}
{}
{}
{}
{}

\bigskip 

\Layout{27}
{}
{}
{}
{./images/wikipedia-cc-by-sa/Sender-Empfaenger-Modell-edited}
{Das Sender-Empfänger-Modell \cite{Informatik3}}
{User:Wiska Bodo, original uploader Freak222
auf \url{http://de.wikipedia.org}, modifiziert; ; CC-BY-SA
3.0}



\Text
{Aussagen einer Nachricht}
{Jede Nachricht hat vier Aussagen\ZFootnote{Nach
Friedemann Schulz von Thun, \cite{ThunRedenI46}}:

\begin{itemize}
  \item \EE{Sachaussage}: Hier vermittelt der
  Sender Daten, Fakten und Sachverhalte.
  \item \EE{Beziehungsaussage}: Hier
  kommt zum Ausdruck, wie der Sender und Empfänger sich
  zueinander verhalten und sich gegenseitig einschätzen.
  Entscheidend ist die Art der Formulierung, die
  Körpersprache, der Tonfall usw., womit Wertschätzung, Respekt,
  Wohlwollen, Gleichgültigkeit, Verachtung etc. vermittelt
  werden kann.
  \item \EE{Appellaussage}: Der Appellanteil einer Nachricht
  soll den Empfänger veranlassen, etwas zu tun
  (oder nicht zu tun). 
  Dieser Versuch einer Einflussnahme kann offen oder
  verdeckt sein, er kann von offenen Bitten und
  Aufforderungen bis hin zu verdeckten
  Manipulationsversuchen reichen.
  \item \EE{Selbstoffenbarung} des Senders: Jede Sendung
  führt unweigerlich zu einer Selbstdarstellung und Selbstenthüllung.
  Sie kann somit zu Deutungen über die
  Persönlichkeit des Sprechers verwendet werden. 
\end{itemize}

Entscheidend in der Kommunikation ist die Tatsache, dass die
gleiche Nachricht beim Sender und beim Empfänger auf
unterschiedliche Art und Weise gedeutet und gewichtet wird,
was zu teils katastrophalen Missverständnissen führen kann. Als
klassisches Beispiel ist die \enquote*{Ampelsituation}
bekannt: 

\begin{quote}
  Ein Paar sitzt im Auto, die Frau sitzt am Steuer, der Mann
  am Beifahrersitz. Sie stehen vor einer roten Ampel, welche
  auf grün wechselt.
  
  Er sagt: \Speech{\enquote{Du, da vorne ist grün!}}
  
  \E{Eine} mögliche Deutung:\small
  
  \begin{tabularx}{\linewidth}{XXX}
  & \TabHeading{Er (Sender)} & \TabHeading{Sie (Empfänger)}
  \\\addlinespace
  \noindent\TabSub{Sachaussage}
  & 
  Die Ampel ist grün.
  &
  Die Ampel ist grün.
  \\\addlinespace
  \noindent\TabSub{Beziehungs\-aus\-sage}
  & 
  Du brauchst meine Hilfe weil Du unaufmerksam bist.
  &
  Er fühlt sich mir überlegen.
  \\\addlinespace
  \noindent\TabSub{Appellaussage}
  &
  Fahr los!
  &
  Er will dass ich ihm das Steuer überlasse.
  \\\addlinespace
  \noindent\TabSub{Selbst\-offen\-barung} 
  &
  Ich habe es eilig.
  &
  Er will mich bevormunden.
  \end{tabularx}
  \normalsize
\end{quote}

}
{}
{}
{}
{}



\Layout{27}
{}
{}
{}
{./images/wikipedia-cc-by-sa/SchulzVonThunVierOhrenModell-edited.pdf}
{Vier-Ohren-Modell der Kommunikation, nach Friedemann Schulz
von Thun \cite{ThunRedenI46}}
{User \enquote*{Mussklprozz} auf Wikipedia, modifiziert;
CC-BY-SA 3.0}




% Der Erfolg des Patienten- bzw. Anamnesegesprächs ist für das
% weitere Vorgehen entscheidend. Oft erhält man wichtige
% Information über den Zustand des Patienten nicht durch die
% gemessenen Werte, sondern durch eine professionelle und
% zielgerichtete Befragung. Neben dem Informationsgewinn wirkt dies zusätzlich
% vertrauensfördernd.
% Der Erfolg des Gespräches hängt dabei einerseits vom Geschick des Gesprächführenden,
% andererseits aber auch von der Bereitschaft des Patienten
% (oder dessen Umgebung!) ab.

\TextSlide{Allgemeine Grundregeln und Ratschläge}{
Ein gut geführtes Patientengespräch ist durch folgende Merkmale gekennzeichnet:
\begin{itemize}
    \item \EE{Vorstellen}: Zur Vorstellung gehört die Nennung des eigenen
    Namens, der Funktion und -- so es die Dringlichkeit zulässt -- ein Händedruck.
    \item \EE{Augenhöhe}: Geht man in die Hocke oder organisiert man sich einen Sessel, 
    kann man mit dem Patient auf
    \E{Augenhöhe} kommunizieren. Dies nimmt dem Patient Angst und fördert eine
    Gesprächsbasis \E{auf einer Ebene}.
    \item \EE{Deutlichkeit und Einfachheit}: Eine deutliche und einfache
    Ausdrucksweise ist besonders wichtig um Missverständnisse zu vermeiden und
    schnell zur gewünschten Information zu kommen.\footnote{Evtl.
    benötigt der Patient Hörhilfen.
    Diese müssen evtl. erst eingeschalten werden.} Fachbegriffe sind zu
    vermeiden!
    \item \EE{Aktives Zuhören}: Die
    Antworten des Patienten müssen bei Helfer \enquote*{ankommen} und nicht
    \enquote*{verloren gehen}. Wichtige Informationen werden oft
    nur einmal berichtet, der Patient nimmt im weiteren
    Verlauf an, dass er sie bereits gesagt und das
    Fachpersonal sie verstanden hat.
    
    Negativbeispiel: \begin{quote}\begin{inparadesc}
                     \item[Sanitäter:] \Speech{\enquote{Haben Sie Schmerzen?}}
                     \item[Patient:] \Speech{\enquote{Nein.}}
                     \item[Sanitäter:] \Speech{\enquote{Wie stark ist der Schmerz?}}
                     \end{inparadesc}\end{quote}
    \item \EE{Verstehen}: Es ist wichtig zu verstehen was der
    Patient mit dem Gesagten \E{meint}! Um sich zu vergewissern, dass man die
    Information richtig verarbeitet hat, kann man das Gesagte in
    eigenen Worten wiederholen und sich vom Patient die Richtigkeit bestätigen
    lassen.
    \item \EE{\enquote*{Monogam}}: Es ist immer nur \E{eine
    Frage gleichzeitig} zu stellen.
    \item \EE{Information}: Der Patient ist immer über das Vorgehen der Helfer
    zu informieren!
    \item \EE{Aufrichtigkeit}: Die Fragen des
    Patienten sind nach bestem Wissen zu beantworten! \EE{Das Anlügen des
    Patienten ist \E{tabu}!} Liegen mögliche Antworten außerhalb des Wissens-
    oder Kompetenzbereiches ist auf die zuständige Stelle zu verweisen.
    \item \EE{Bezugsperson}: Sind mehrere Helfer
    vorort, so soll nur einer der Ansprechpartner des
    Patienten sein!\label{topic:einer-redet}
\end{itemize}
}{
\begin{itemize}
  \item Vorstellen
  \item auf Augenhöhe mit dem Patienten
  \item deutliche, einfache Sprache
  \item aktives Zuhören
  \item Verstehen der Antworten
  \item nur eine Frage gleichzeitig
  \item Informieren über die Maßnahmen
  \item Aufrichtigkeit
  \item Bezugsperson
\end{itemize}
}{}{}{}



\TextSlide{Fragentypen}
{
Grundsätzlich unterscheidet man \EE{offene} und
\EE{geschlossene}, sowie \EE{gezielte} Fragen und
\EE{Suggestivfragen}.

\begin{itemize}
  \item \EE{Offene Fragen} kann der Patient frei, mit
  seinen eignen Worten beantworten. Je nach Patient kann man so
  mit wenig Aufwand sehr viele Informationen gewinnen. Das
  Problem dabei kann jedoch sein, dass die Informationen
  ungeordnet und mitunter auch in der Situation mäßig wichtig
  sind.
  
  \begin{quote}
    \Speech{\enquote{Was haben Sie für ein Problem?}}
    
    \Speech{\enquote{Bitte beschreiben Sie den Schmerz.}}
  \end{quote}
  
  
  \item \EE{Geschlossene Fragen} sind sehr direkt und
  speziell, der Patient kann i.\,d.\,R. nur mit
  \emph{\enquote{ja}}, \emph{\enquote{nein}} oder
  \emph{\enquote{weiß nicht}} antworten. Sie machen daher
  eher nur dann Sinn, wenn man ganz spezielle Dinge bewusst
  direkt abfragen kann oder muss. \Ca{Achtung!} Manche
  Patienten neigen dazu, grundsätzlich \enquote{ja} zu
  sagen, ohne dass sie die Frage wirklich verstanden haben!
  
  \begin{quote}
    \Speech{\enquote{Haben Sie sich vorher angestrengt?}}
    
    \Speech{\enquote{Haben Sie das schon einmal gehabt?}}
  \end{quote}
  
  \item \EE{Gezielte Fragen} fokussieren die Antwort des
  Patienten und
  lassen sich nur schwer \enquote{umgehen}. Sie dienen der
  Präzisierung. Werden zu
  viele von diesen Fragen gestellt, kann das Gespräch leicht
  in eine \enquote{Einbahnstraße} abgleiten. Gezielte Fragen können
  sowohl offen, als auch geschlossen sein.
  \begin{quote}
    \Speech{\enquote{Wohin strahlt der
    Schmerz aus?}}
  \end{quote}
  
  \item \EE{Suggestivfragen} sind \enquote*{böse}: Hier
  legt der Gesprächsführende dem Patienten die Antworten \enquote*{in
  den Mund}. Sie können sowohl als offene, als auch als
  geschlossene Fragen formuliert werden.
  
  \begin{quote}
    \Speech{\enquote{Haben Sie \emph{eh keine} Schmerzen?}}
    
    \Speech{\enquote{Ist der Schmerz stechend?}}
  \end{quote}
  
  Dieser Fragentyp soll nicht -- oder nur im begründeten
  Ausnahmefall -- angewendet werden.
\end{itemize}



}
{
\begin{itemize}
    \item Offene Fragen
    \item Geschlossene Fragen
    \item Gezielte Fragen
    \item Suggestivfragen 
\end{itemize}
}{}{}{}

\TextSlide{Kinder}
{Die Kommunikation mit Kindern kann besonders herausfordernd
sein. Der Kontakt mit Einrichtungen des Gesundheitswesens,
sei es in einer Ordination, im Spital oder im Rahmen des
Rettungsdienstes, stellt eine aufregende und i.\,d.\,R.
angsteinflößende Ausnahmesituation dar. Die folgenden
Vorschläge können unter Umständen hilfreich sein -- ein
allgemeingültiges \enquote{Kochrezept} gibt es allerdings nicht:

\begin{itemize}
  \item \EE{Kleine Erwachsene}: Manchmal sind Kinder doch
  einfach \enquote{kleine Erwachsene}: Die elementaren Grundsätze
  \E{Aufrichtigkeit} und \E{Aufklärung}, welche für das
  Verhalten gegenüber Erwachsenen gilt, gilt genauso für
  Kinder. Ein altersangepasstes Auftreten ist zwar oft
  notwendig, darf aber nicht als Ausrede benutzt werden, um
  die Grundsätze zu vernachlässigen.
%   \item \EE{Aufrichtigkeit}: Seien Sie immer ehrlich! Kinder
%   merken sehr schnell, wenn Sie lügen. Lügen  Sie ein Kind
%   an, so wird das Kind schnell das Vertrauen in Sie
%   verlieren  und jede weitere Untersuchung verweigern!
% 
%   Kinder nehmen non-verbale Signale stärker wahr als
%   Erwachsene. Versuchen Sie daher durch einen freundlichen
%   Gesichtsausdruck und eine offene Körperhaltung Vertrauen
%   zu bilden!
  \item \EE{Aufrichtigkeit}: Kinder bemerken sehr schnell, wenn man lügt. Ein
  ehrlicher Umgang mit Kindern ist daher besonders wichtig, da sonst das
  Vertrauen verloren ist und das Kind möglicherweise spätere Untersuchungen
  (aus Angst) verweigert.
%   \item \EE{Erklären} Sie dem Kind (altersgerecht) warum Sie
%   mit ihm oder ihr reden oder es untersuchen möchten.
  \item \EE{Erklären}: Das (altersgerechte) Erklären von Untersuchungen und
  Maßnahmen kann dem Kind Angst nehmen.
%   \item \EE{Reden} Sie mit dem Kind und sprechen sie
%   es direkt und mit seinem Namen an! Vermeiden  Sie es über
%   das Kind in der 3. Person zu sprechen, damit das Kind
%   nicht das Gefühl  bekommt, dass hinter seinem Rücken
%   getuschelt wird. Dies würde dem Kind unnötig Angst machen!
  \item \EE{Direktes Ansprechen}: Kinder sind immer direkt anzusprechen!
  Unterhält man sich ständig mit den Eltern über das Kind, so kann es das Gefühl
  haben, dass hinter seinem Rücken getuschelt wird. Dies würde dem Kind unnötig
  Angst machen!
%   \item Geben Sie dem Kind ein Spielzeug in die Hand!
  \item \EE{Spielzeug} kann Kinder (von der Angst) ablenken. 
%   \item \EE{Bezugsperson}: Trennen Sie das Kind niemals von
%   der Bezugsperson! Hat das Kind Angst vor einer bestimmten
%   Untersuchung oder einer Maßnahme, zeigen Sie diese zuerst
%   an der Bezugsperson oder an einem Stofftier vor!
  \item \EE{Bezugsperson}: Kinder brauchen die Nähe ihrer Eltern
  (Bezugspersonen)! Kinder daher niemals von diesen trennen. Außerdem kann man
  bevorstehende Untersuchungen oder Maßnahmen an Bezugspersonen (oder
  Stofftieren) vorzeigen, sodass dem Kind Angst genommen wird.
%  \item Loben Sie das Kind nach der Untersuchung bzw. nach der Maßnahme!
  \item \EE{Lob} hilft ebenfalls weiter!
%   \item \EE{Zeigen lassen}: Da Kinder ihre Beschwerden oft nicht
%   eindeutig beschreiben können, lassen Sie sich z.\,B. den
%   Ort der Schmerzen zeigen. Auf die \E{Körperhaltung} u.\,ä.
%   äußere Anzeichen muss besonders
%   geachtet werden!
  \item \EE{Zeigen lassen}: Da Kinder ihre Beschwerden oft nicht
  eindeutig beschreiben können, ist es sinnvoll sich z.\,B. den
  Ort der Schmerzen zeigen zu lassen. Außerdem muss besonders auf die
  \E{Körperhaltung} u.\,ä.
  äußere Anzeichen geachtet werden!
\end{itemize}
}{
\begin{itemize}
  \item \E{Aufrichtigkeit} und \E{Aufklärung} gelten auch
  für Kinder
  \item Aufrichtigkeit
  \item Aufklärung
  \item Erklären
  \item Reden
  \item Bezugsperson
  \item Sachen \E{zeigen} lassen
\end{itemize}
}{}{}{}
